\documentclass[11pt,a4paper]{article}
\usepackage[T1]{fontenc}
\usepackage{lmodern,amsmath,amssymb,amsthm,mathtools,booktabs,microtype}
\usepackage[margin=27mm]{geometry}
\usepackage[colorlinks=true,linkcolor=blue,citecolor=blue,urlcolor=blue]{hyperref}
\setlength{\emergencystretch}{2em}
\newtheorem{theorem}{Theorem}[section]
\newtheorem{lemma}[theorem]{Lemma}
\newtheorem{proposition}[theorem]{Proposition}
\newtheorem{corollary}[theorem]{Corollary}
\newcommand{\Q}{\mathbb Q}
\DeclareMathOperator{\Tr}{Tr}\DeclareMathOperator{\Nm}{N}
\title{High-genus fixed fibers of a native square criterion}
\author{}\date{16 November 2026}
\begin{document}\maketitle
\begin{abstract}
We compute the genera of four fixed fibers of the rational-square
criterion for the faster conditional elliptic detector. The input
values are one and two, with actual source points
$A=[X]P$, $H=3P$, and $U=[3X]P$, and with coefficient shift three
and height exponent four. Both the original and symmetric final
comparison are treated. The single square covers have genera $59$
and $83$, and their geometrically integral simultaneous-cover fibers
have genera $236$ and $332$. Exact good-reduction certificates prove
the characteristic-zero squarefreeness and branch intersections.
Each simultaneous fiber has four rational points at infinity, which
are outside the guarded source domain. No nonexistence of finite
rational points, obstruction to auxiliary-index variation, or smaller
defining-variable bound is asserted.
\end{abstract}

\section{Actual point data and detector functions}
Put
\[
 K=\Q(\sqrt2),\qquad d=3-\sqrt2,\qquad
 E_K:\ y^2=x^3-d^2x,
\]
and use the specified source point
\[
 P=\left(\frac{57-28\sqrt2}{8},
          \frac{-660+417\sqrt2}{32}\right)
   =2(1+\sqrt2,2\sqrt2).
\]
For each $X\in\{1,2\}$ fix
\begin{equation}\label{eq:points}
                  A=[X]P,\qquad H=3P,\qquad U=[3X]P.
\end{equation}
Thus $H-P=2P$ and $U-A=[2X]P$. These are actual doubled source
points, rather than free coordinates on a norm conic. All ordinates,
abscissas and difference denominators used below are nonzero; the
exact certificate verifies this on the displayed points.

Write $f_0(Q)=1/x(Q)$, $z(Q)=-x(Q)/y(Q)$ and
\[
 r=\frac{z(U)}{z(H)},\qquad
 s=\frac{z(U-A)}{z(H-P)}.
\]
Conjugation in $K/\Q$ is denoted by a bar. The first five tests are
\[
 C_1=f_0(H),\quad C_2=f_0(U),\quad
 C_3=\overline{f_0(H-P)},\quad C_4=\overline{f_0(U-A)},
 \quad C_5=r-\bar s.
\]
We consider both choices
\begin{equation}\label{eq:last}
       C_6^{\mathrm{orig}}=r-X,
       \qquad C_6^{\mathrm{sym}}=r+\bar s-2X.
\end{equation}
They are the original conditional comparison of~\cite{Direct} and
the symmetric replacement of~\cite{Criterion}, respectively.
All six tests are fixed elements of $K$ once~\eqref{eq:points} is
fixed.

Let
\begin{equation}\label{eq:F}
 F(t)=\frac{t(t+1)(2t+1)(3t+1)(4t+1)(5t+2)}{120},\qquad
 \delta=t-1,\qquad w(t)=\frac{F(t)}{4\delta^8}.
\end{equation}
Use the faster quadratic product
\begin{equation}\label{eq:product}
                    W(t)=\prod_{i=1}^{6}\left(1+\frac{C_i^2}{F(t)}\right).
\end{equation}
Set $V=\Tr(W^{-1})$, $J=\Nm(W^{-1})$, $g=1+w$, and
\begin{equation}\label{eq:criterion}
 Q=1-V/g+J/g^2,\qquad
 \mathcal R(t)=Q\left(1+\frac{1-V/2}{w}\right).
\end{equation}
This is the flat member of the finite square criterion
of~\cite{Criterion}. All functions retain their source denominators
and guards. The complete curves below mean the smooth projective
models of the indicated function fields; the source domain is an
open subset of those models.

For the coefficient shift $B=3$ and height exponent four, the two
actual arguments are $t_0=\lambda$ and $t_1=a\lambda$, where
\begin{equation}\label{eq:scale}
                     a=(X^2+3)^4+\bigl(\Tr x(A)\bigr)^4.
\end{equation}
The value $a$ is the source scale, not an independently selected
rational coefficient.

\section{A polynomial model of the square cover}
Let $m$ be the number of nonzero tests among $C_1,\ldots,C_6$.
For the two conjugate embeddings write
\[
 P_\pm(t)=\prod_{C_i\ne0}\bigl(F(t)+C_{i,\pm}^2\bigr),\qquad
 D_0(t)=F(t)^m,\qquad W_\pm=P_\pm/D_0.
\]
The two polynomials $P_\pm$ are conjugate and monic of degree $6m$.
Define
\begin{equation}\label{eq:branchpolynomials}
\begin{gathered}
 E_0(t)=4(t-1)^8,\qquad G(t)=E_0(t)+F(t),\\
 L_\pm(t)=G(t)P_\pm(t)-E_0(t)D_0(t),\\
 M(t)=2G(t)P_+(t)P_-(t)
             -E_0(t)D_0(t)\bigl(P_+(t)+P_-(t)\bigr),\\
                       \mathfrak f(t)=2F(t)L_+(t)L_-(t)M(t).
\end{gathered}
\end{equation}
In particular $\mathfrak f\in\Q[t]$.

\begin{lemma}\label{lem:identity}
On the guarded rational-function domain,
\begin{equation}\label{eq:coveridentity}
         \mathcal R(t)=
          \frac{\mathfrak f(t)}
               {\bigl(2F(t)G(t)P_+(t)P_-(t)\bigr)^2}.
\end{equation}
Moreover
\[
 \deg\mathfrak f=24m+24,\qquad
                  \text{the leading coefficient of }\mathfrak f\text{ is }4.
\]
\end{lemma}
\begin{proof}
Let $N=\Nm(W)$ and $A_0=\Tr(W)/2$. The rational identity for the
flat criterion is
\[
 \mathcal R=
 \frac{(g^2N-2gA_0+1)(gN-A_0)}{wg^2N^2}.
\]
Substitute $N=P_+P_-/D_0^2$, $2A_0=(P_++P_-)/D_0$,
$w=F/E_0$, and $g=G/E_0$. The result is
\[
 \mathcal R=
 \frac{L_+L_-M}{2FG^2(P_+P_-)^2},
\]
which is~\eqref{eq:coveridentity}.

Since $P_\pm-D_0$ has degree at most $6m-6$,
\[
 L_\pm=E_0(P_\pm-D_0)+FP_\pm
\]
has degree $6m+6$ and leading coefficient one. Similarly,
\[
 M=2FP_+P_-+E_0\bigl((P_+-D_0)P_-+(P_--D_0)P_+\bigr)
\]
has degree $12m+6$ and leading coefficient two. Multiplying in
\eqref{eq:branchpolynomials} proves the degree and coefficient
assertions.
\end{proof}

Thus the rational square test has the same quadratic function-field
extension as $y^2=\mathfrak f(t)$. This statement does not remove
any of the source guards.

\section{Squarefreeness and the simultaneous-cover genus}

\begin{theorem}\label{thm:genera}
For the actual point tuples~\eqref{eq:points}, with either final
comparison in~\eqref{eq:last}, the polynomial $\mathfrak f$ is
squarefree. Its degree is $120$ at $X=1$ and $168$ at $X=2$.
The corresponding single covers have genera $59$ and $83$.

For the actual scale~\eqref{eq:scale},
\[
            \gcd\bigl(\mathfrak f(t),\mathfrak f(at)\bigr)=t
\]
up to multiplication by a nonzero constant. The simultaneous-cover
function field
\[
 \Q(t)\left(\sqrt{\mathfrak f(t)},\sqrt{\mathfrak f(at)}\right)
\]
is geometrically integral of degree four over $\Q(t)$. Its smooth
projective curve has genus $236$ at $X=1$ and $332$ at $X=2$.
\end{theorem}

\begin{proof}
At $X=1$ one has $A=P$ and $U=H$, so $r=s=1$. Both ratio tests
vanish and exactly four reciprocal tests remain: $m=4$. At $X=2$
all six tests are nonzero, as certified by their nonzero reductions
in the table below. Lemma~\ref{lem:identity} gives the two degrees.

The exact good-reduction data are as follows. Each prime is only a
verification prime; it need not lie outside the detector's arithmetic
exceptional set.
\begin{center}
\begin{tabular}{llrrrr}\toprule
$X$ & Final test & Prime & $a\bmod p$ & Degree & Branch gcd\\\midrule
1 & Original & 97 & 71 & 120 & $t$\\
1 & Symmetric & 97 & 71 & 120 & $t$\\
2 & Original & 97 & 79 & 168 & $t$\\
2 & Symmetric & 241 & 125 & 168 & $t$\\\bottomrule
\end{tabular}
\end{center}
All rational coefficient denominators, actual point ordinates and
abscissas, and affine difference denominators used in constructing
these reductions are units at the indicated prime. The certificate
records the actual $K$-coordinates, their reductions, the complete
coefficient vectors in~\eqref{eq:branchpolynomials}, and the two
explicit polynomial identities
\begin{equation}\label{eq:bezout}
\begin{gathered}
 s_0\mathfrak f+t_0\mathfrak f'=1,\\
 s_1\bigl(\mathfrak f/t\bigr)
       +t_1\bigl(\mathfrak f(at)/t\bigr)=1
                         \qquad\text{in }\mathbb F_p[t].
\end{gathered}
\end{equation}
These identities are verified by the independent standard-library
checker using polynomial addition and multiplication, without a
factorization package. Both reduced polynomials retain their full
degrees. A characteristic-zero common factor of positive degree
would retain that degree after monic normalization at this prime:
it divides a polynomial with unit leading coefficient and integral
coefficients, so its coefficients are integral. Therefore the first
identity proves squarefreeness over $\Q$.

The root zero is simple in $F$. At zero every nonzero test is still
nonzero, $P_\pm(0)\ne0$, and the expressions for $L_\pm,M$ are
nonzero. Hence $t$ divides $\mathfrak f$ exactly once. Since
$a\ne0$, it also divides $\mathfrak f(at)$ exactly once. The second
identity in~\eqref{eq:bezout} excludes every further
characteristic-zero common factor.

Write $n=\deg\mathfrak f$. The single cover has $n$ simple finite
branch points and is unramified at infinity because $n$ is even.
Its genus is $(n-2)/2$ by Riemann--Hurwitz~\cite{RH}.
The two branch sets have only zero in common. Their square classes
are independent over $\overline\Q(t)$: either polynomial has a
simple root absent from the other, and their product is also
nonsquare. The simultaneous cover is consequently geometrically
integral of degree four. Its branch-set union has size $2n-1$.
Above each branch point there are two points of ramification index
two, including at the common branch point. Riemann--Hurwitz gives
\[
                 2g-2=-8+2(2n-1),\qquad g=2n-4,
\]
which gives the stated genera.
\end{proof}

The intermediate identities are supplied as coefficient vectors rather
than only as hashes or factorization reports. Recomputing the source
curve operations and all modular polynomial products suffices to audit
the characteristic-zero conclusions.

\section{Rational boundary points and the unresolved interior}

\begin{proposition}\label{prop:infinity}
Each simultaneous-cover curve in Theorem~\ref{thm:genera} has exactly
four rational points above $t=\infty$.
\end{proposition}
\begin{proof}
If $n=\deg\mathfrak f$, the local coordinate $y/t^{n/2}$ on the
first quadratic cover has limiting square equal to four. Thus the
two points at infinity are rational, with limiting values $2,-2$.
For the second cover the corresponding values are
$2a^{n/2},-2a^{n/2}$, again rational because $a\in\Q^\times$.
Both covers are unramified there, so their four ordered pairs give
exactly four rational points on the simultaneous cover.
\end{proof}

The parameter infinity is not in the guarded source domain. These
boundary points therefore give no detector witness at a finite common
slope. In particular, the geometrically integral curves in
Theorem~\ref{thm:genera} have no rational or elliptic component on
which a chord or tangent construction from the boundary would provide
the missing witnesses.

The theorem does not assert that the curves have no finite rational
points. It treats only the four displayed fixed fibers, with $B=3$
and exponent four. It does not restrict another choice of auxiliary
indices, a proper coupled family of actual points, or a different
square criterion. The simultaneous actual rational-square property
needed for the twenty-one- and nine-variable constructions remains
unresolved.

\section{Verification package}
The accompanying files are \nolinkurl{certificate.json},
\nolinkurl{generate.py}, and \nolinkurl{check.py}. Generation uses
exact rational arithmetic in $\Q(\sqrt2)$ and modular polynomial
arithmetic. The independent checker uses only the Python standard
library and verifies the complete coefficient and B\'ezout data,
the cyclic source-point reductions, the actual scale, and the source
file hashes. It does not search for rational points on these curves
or replace a global realization theorem by local evidence.

\begin{thebibliography}{9}
\bibitem{Direct}\emph{Thirteen variables from direct elliptic index
comparison}, repository manuscript,
\nolinkurl{papers/conditional_few_variables/direct_input13_20261110_v22.tex}.
\bibitem{Criterion}\emph{A guarded rational-square criterion for
normalized elliptic detectors}, companion manuscript, 16 November 2026.
\bibitem{RH} The Stacks Project Authors, \emph{Riemann--Hurwitz},
Section 53.12, Tag 0C1B,
\url{https://stacks.math.columbia.edu/tag/0C1B}.
\end{thebibliography}
\end{document}
